\documentclass[a4paper,11pt]{article}

% ============================================================
%  Le Collé — Algèbre linéaire et réduction
%  © 2026 Ahmed Amine Ouqlani
% ============================================================

% ---------- Polices ----------
\usepackage[utf8]{inputenc}
\usepackage[T1]{fontenc}
\usepackage[sc]{mathpazo}          % texte et mathématiques : Palatino
\usepackage[scaled=0.88]{helvet}   % titres et repères : linéale
\linespread{1.07}
\usepackage{amsmath,amssymb,amsthm,mathrsfs,stmaryrd}
\usepackage{microtype}

% ---------- Mise en page ----------
\usepackage[a4paper,top=2.3cm,bottom=2.3cm,left=2.4cm,right=2.4cm,
            headheight=14pt,headsep=0.8cm,footskip=1.1cm]{geometry}
\setlength{\parindent}{0pt}
\setlength{\parskip}{0.45em plus 0.15em minus 0.1em}
\raggedbottom
\emergencystretch=1.5em
\clubpenalty=10000 \widowpenalty=10000 \displaywidowpenalty=10000 \brokenpenalty=10000

% ---------- Couleurs : un bleu de nuit, un ocre, des gris ----------
\usepackage{xcolor}
\definecolor{Ink}{HTML}{1C2533}
\definecolor{Bleu}{HTML}{1F3A5F}
\definecolor{BleuMoyen}{HTML}{4A6A8F}
\definecolor{BleuClair}{HTML}{9DB2C9}
\definecolor{Fond}{HTML}{F2F4F7}
\definecolor{Ocre}{HTML}{B07A26}
\definecolor{Brique}{HTML}{A23B2A}
\definecolor{Gris}{HTML}{5E6672}
\definecolor{Filet}{HTML}{C5CCD6}

% ---------- Packages ----------
\usepackage{enumitem}
\setlist{itemsep=0.25em,topsep=0.35em,parsep=0pt}
\setlist[itemize]{label=\textcolor{BleuMoyen}{\small\textbullet},leftmargin=1.4em}
\setlist[enumerate]{leftmargin=1.8em,label=\textsf{\small\color{Bleu}\arabic*.}}
\usepackage{booktabs,array,multicol}
\usepackage{needspace,float}
\usepackage[font=small,labelfont={sf,color=Bleu},labelsep=period,skip=6pt]{caption}
\usepackage{tikz}
\usetikzlibrary{arrows.meta,calc,decorations.markings}
\usepackage{pgfplots}
\pgfplotsset{compat=1.17}
\usepackage{tkz-tab}
\usepackage[most]{tcolorbox}
\usepackage{titlesec,titletoc}
\usepackage{fancyhdr,lastpage,etoolbox}
\usepackage[colorlinks=true,linkcolor=Bleu,urlcolor=Bleu,citecolor=Bleu,
            pdftitle={Le Collé — Algèbre linéaire et réduction},pdfauthor={Ahmed Amine Ouqlani}]{hyperref}

\numberwithin{equation}{section}
\numberwithin{figure}{section}
\numberwithin{table}{section}

% ---------- Titres ----------
% Chapitres (sections) : grand numéro, titre en linéale, filet ocre
\titleformat{\section}[display]
  {\normalfont\sffamily\color{Ink}}
  {{\small\color{Gris}\textls[200]{\MakeUppercase{Chapitre}}}\\[-0.2em]%
   {\fontsize{64}{64}\selectfont\color{BleuClair}\thesection}}
  {0.4em}
  {\fontsize{25}{30}\selectfont\bfseries}
  [\vspace{0.5em}{\color{Ocre}\rule{3.2cm}{1.4pt}}]
\titlespacing*{\section}{0pt}{0pt}{1.6em}
\titleformat{name=\section,numberless}[display]
  {\normalfont\sffamily\color{Ink}}{}{0pt}{\fontsize{25}{30}\selectfont\bfseries}
  [\vspace{0.5em}{\color{Ocre}\rule{3.2cm}{1.4pt}}]
\titlespacing*{name=\section,numberless}{0pt}{0pt}{1.6em}

\titleformat{\subsection}
  {\needspace{5\baselineskip}\normalfont\sffamily\Large\bfseries\color{Ink}}
  {\color{Bleu}\thesubsection}{0.8em}{}
\titlespacing*{\subsection}{0pt}{1.8em}{0.5em}
\titleformat{\subsubsection}
  {\needspace{5\baselineskip}\normalfont\sffamily\large\bfseries\color{Ink}}
  {\color{Bleu}\thesubsubsection}{0.7em}{}
\titlespacing*{\subsubsection}{0pt}{1.6em}{0.4em}
\titleformat{\paragraph}[runin]{\normalfont\sffamily\bfseries\color{Ink}}{}{0pt}{}[.]
\titlespacing*{\paragraph}{0pt}{0.9em}{0.6em}
\setcounter{secnumdepth}{3}
\setcounter{tocdepth}{2}

% ---------- Table des matières ----------
\titlecontents{section}[0pt]{\addvspace{0.7em}\sffamily\bfseries\color{Ink}}
  {\makebox[2.2em][l]{\color{Bleu}\thecontentslabel}}{}{\hfill\contentspage}[\addvspace{0.2em}]
\titlecontents{subsection}[2.2em]{\small\linespread{1}\selectfont}
  {\makebox[2.6em][l]{\color{Gris}\thecontentslabel}}{}{\titlerule*[0.7pc]{.}\contentspage}

% ---------- En-têtes et pieds de page ----------
\pagestyle{fancy}
\fancyhf{}
\renewcommand{\sectionmark}[1]{\markboth{\thesection\quad#1}{}}
\fancyhead[L]{\footnotesize\sffamily\color{Gris}\leftmark}
\fancyhead[R]{\footnotesize\sffamily\scshape\color{Gris}Le Collé}
\fancyfoot[C]{\small\sffamily\color{Ink}\thepage}
\fancyfoot[L]{\scriptsize\sffamily\hyperlink{toc}{\color{Filet}sommaire}}
\renewcommand{\headrulewidth}{0.4pt}
\renewcommand{\headrule}{{\color{Filet}\hrule height\headrulewidth width\headwidth}}
\renewcommand{\footrulewidth}{0pt}
\fancypagestyle{plain}{\fancyhf{}\fancyfoot[C]{\small\sffamily\color{Ink}\thepage}\renewcommand{\headrulewidth}{0pt}}

% ---------- Énoncés ----------
\newtheoremstyle{enonce}{0pt}{0pt}{}{}{\sffamily\bfseries\color{Bleu}}{.}{0.5em}
  {\thmname{#1}\thmnumber{~#2}\thmnote{\ \normalfont\sffamily\color{Ink}(#3)}}
\newtheoremstyle{exemple}{0.9em}{0.9em}{}{}{\sffamily\bfseries\color{BleuMoyen}}{.}{0.5em}
  {\thmname{#1}\thmnumber{~#2}\thmnote{\ \normalfont\sffamily\color{Ink}(#3)}}
\newtheoremstyle{remarque}{0.7em}{0.7em}{}{}{\sffamily\itshape}{.}{0.5em}{\thmname{#1}\thmnote{\ (#3)}}
\theoremstyle{enonce}
\newtheorem{definition}{Définition}[section]
\newtheorem{theoreme}[definition]{Théorème}
\newtheorem{proposition}[definition]{Proposition}
\newtheorem{corollaire}[definition]{Corollaire}
\theoremstyle{exemple}
\newtheorem{exemple}[definition]{Exemple}
\theoremstyle{remarque}
\newtheorem*{remarque}{Remarque}

\tcbset{
  defi/.style={enhanced,breakable,colback=white,colframe=white,boxrule=0pt,arc=0pt,
     borderline west={1.2pt}{0pt}{Bleu},left=11pt,right=4pt,top=3pt,bottom=3pt,
     before skip=1.1em,after skip=1.1em},
  thm/.style={enhanced,breakable,colback=Fond,colframe=Fond,boxrule=0pt,arc=0pt,
     borderline west={2.4pt}{0pt}{Bleu},left=11pt,right=8pt,top=6pt,bottom=6pt,
     before skip=1.2em,after skip=1.2em}}
\tcolorboxenvironment{definition}{defi}
\tcolorboxenvironment{theoreme}{thm}
\tcolorboxenvironment{proposition}{thm}
\tcolorboxenvironment{corollaire}{thm}

\renewcommand{\proofname}{Démonstration}
\renewcommand{\qedsymbol}{\textcolor{BleuClair}{$\blacksquare$}}
\makeatletter
\renewenvironment{proof}[1][\proofname]{\par\pushQED{\qed}\normalfont
  \topsep0.5em\relax\trivlist\item[\hskip\labelsep\sffamily\itshape\color{Gris}#1.]\ignorespaces}
  {\popQED\endtrivlist\@endpefalse}
\makeatother

% ---------- Méthode : filets haut et bas, titre en bandeau ----------
\makeatletter\newcommand{\methlabelfix}{\protected@edef\@currentlabel{\thetcbcounter}}\makeatother
\newtcolorbox[auto counter,number within=section]{methode}[2][]{%
  enhanced,unbreakable,colback=white,colframe=white,boxrule=0pt,arc=0pt,
  borderline north={0.9pt}{0pt}{Bleu},borderline south={0.9pt}{0pt}{Bleu},
  left=4pt,right=4pt,top=8pt,bottom=8pt,before skip=1.4em,after skip=1.4em,
  before upper={\methlabelfix\phantomsection\addcontentsline{met}{metentry}{\protect\makebox[2.6em][l]{\thetcbcounter}#2}%
     {\sffamily\bfseries\small\color{white}\colorbox{Bleu}{\strut\ MÉTHODE~\thetcbcounter\ }}%
     \ \ {\sffamily\bfseries\color{Bleu}#2}\par\smallskip},#1}

% ---------- Repères pédagogiques : un filet de couleur, un mot-clé ----------
\newcommand{\definirbloc}[3]{% nom, couleur, libellé
  \newtcolorbox{#1}[1][]{enhanced,breakable,blanker,
    borderline west={1.4pt}{0pt}{#2},left=10pt,right=2pt,top=1pt,bottom=1pt,
    before skip=1em,after skip=1em,
    before upper={{\sffamily\bfseries\footnotesize\color{#2}\textls[120]{\MakeUppercase{#3}}}\quad},##1}}
\definirbloc{reflexe}{Bleu}{Réflexe}
\definirbloc{astuce}{Ocre}{Astuce}
\definirbloc{attention}{Brique}{Attention}
\definirbloc{enpratique}{BleuMoyen}{En pratique}
\definirbloc{plusloin}{Gris}{Pour aller plus loin}
\definirbloc{strategie}{Bleu}{Stratégie}
\definirbloc{lecture}{BleuMoyen}{Lire le schéma}
% compatibilité avec les parties reprises
\newenvironment{garde}{\begin{attention}}{\end{attention}}
\newenvironment{complementHP}{\begin{plusloin}}{\end{plusloin}}
\newenvironment{interpretation}{\begin{enpratique}}{\end{enpratique}}

% ---------- Exercices, corrigés ----------
\newcounter{exo}[section]
\renewcommand{\theexo}{\thesection.\arabic{exo}}
\newcommand{\niveau}[1]{{\color{Ocre}\footnotesize
  \ifcase#1\or$\bullet\circ\circ\circ\circ$\or$\bullet\bullet\circ\circ\circ$\or$\bullet\bullet\bullet\circ\circ$%
  \or$\bullet\bullet\bullet\bullet\circ$\or$\bullet\bullet\bullet\bullet\bullet$\fi}}
\newcommand{\serie}[1]{\par\needspace{8\baselineskip}\vspace{1.5em}%
  \noindent{\sffamily\bfseries\color{Bleu}#1}\par\nobreak\vspace{-0.4em}%
  \noindent{\color{Filet}\rule{\linewidth}{0.4pt}}\par\nobreak\vspace{0.1em}}
\newenvironment{exo}[3]{%
  \par\needspace{5\baselineskip}\vspace{0.75em}\refstepcounter{exo}\label{e:#1}%
  \def\exolab{#1}%
  \noindent{\sffamily\bfseries\color{Bleu}Exercice~\theexo}\quad\niveau{#2}%
  \hfill{\footnotesize\sffamily\color{Gris}#3}\par\nobreak\vspace{0.15em}}%
  {\par\nopagebreak\noindent\hfill{\scriptsize\sffamily\hyperref[c:\exolab]{\color{Gris}corrigé p.~\pageref{c:\exolab}}}\par}
\newenvironment{corrige}[1]{%
  \par\needspace{6\baselineskip}\vspace{1em}\phantomsection\label{c:#1}%
  \noindent{\sffamily\bfseries\color{Bleu}Exercice~\ref{e:#1}}\quad
  {\scriptsize\sffamily\hyperref[e:#1]{\color{Gris}énoncé p.~\pageref{e:#1}}}\par\nobreak\vspace{0.2em}}%
  {\hfill\textcolor{BleuClair}{$\square$}\par}
\newcommand{\etape}[1]{\par\smallskip\noindent{\sffamily\bfseries\small\color{BleuMoyen}#1.}\ }
\newenvironment{solution}{\par\noindent{\sffamily\bfseries\small\color{BleuMoyen}Solution.}\ }{\hfill\textcolor{BleuClair}{$\square$}\par}
\newcommand{\chercher}{\par\smallskip\noindent{\color{Filet}\rule[0.5ex]{1.2cm}{0.4pt}}\
  {\footnotesize\sffamily\itshape\color{Gris}laisser chercher l'étudiant avant de poursuivre}\par\smallskip}

% ---------- Guide de séance ----------

\newcommand{\marge}[1]{\leavevmode\marginpar{\raggedleft\scriptsize\sffamily\color{Gris}\textls[100]{\MakeUppercase{#1}}}\ignorespaces}
\newcommand{\Etape}[4]{%
  \subsection[Étape #1 --- #3 {\color{Gris}(#2)}]{#3}\label{etape:#1}%
  \noindent{\small\sffamily\color{Gris}{\color{Ocre}\bfseries Étape #1}\quad·\quad #2\quad·\quad #4}\par\medskip}
\newcommand{\question}[1]{\par\smallskip\noindent{\sffamily\bfseries\small\color{BleuMoyen}Question.}\ \textit{#1}\par\smallskip}
\newcommand{\rappel}[1]{\par\smallskip\noindent\phantomsection\addcontentsline{rap}{rapentry}{#1}%
  {\sffamily\bfseries\scriptsize\colorbox{Fond}{\color{Gris}\strut RAPPEL}}\ {\sffamily\small\color{Gris}#1.}\ }
\newcommand{\pourquoi}{\par\noindent{\sffamily\bfseries\small\color{Ocre}Pourquoi ?}\ }

% Ouverture de chapitre : \ouverture{objectifs}{prérequis}{durée}
\newcommand{\ouverture}[3]{%
  \begin{tcolorbox}[enhanced,colback=Fond,colframe=Fond,arc=0pt,boxrule=0pt,left=12pt,right=12pt,top=10pt,bottom=10pt,
     before skip=1.2em,after skip=1.6em]
  \begin{minipage}[t]{0.56\linewidth}
  {\sffamily\bfseries\small\color{Bleu}OBJECTIFS}\par\vspace{0.3em}\small #1
  \end{minipage}\hfill
  \begin{minipage}[t]{0.38\linewidth}
  {\sffamily\bfseries\small\color{Bleu}PRÉREQUIS}\par\vspace{0.3em}\small #2
  \ifx\relax#3\relax\else\par\vspace{0.6em}{\sffamily\bfseries\small\color{Bleu}DURÉE}\par\vspace{0.2em}\small #3\fi
  \end{minipage}
  \end{tcolorbox}}
% Fin de chapitre
\newenvironment{essentiel}{\par\needspace{12\baselineskip}\vspace{1.5em}%
  \begin{tcolorbox}[enhanced,unbreakable,colback=white,colframe=white,boxrule=0pt,arc=0pt,
   borderline north={1.4pt}{0pt}{Ocre},left=2pt,right=2pt,top=8pt,bottom=4pt]
  {\sffamily\bfseries\color{Ocre}L'ESSENTIEL DU CHAPITRE}\par\medskip}{\end{tcolorbox}}

% ---------- Notations ----------
\newcommand{\K}{\mathbb{K}}
\newcommand{\R}{\mathbb{R}}
\newcommand{\C}{\mathbb{C}}
\newcommand{\N}{\mathbb{N}}
\newcommand{\Z}{\mathbb{Z}}
\newcommand{\Mn}[1]{\mathcal{M}_{n}(#1)}
\newcommand{\GLn}[1]{\mathrm{GL}_n(#1)}
\newcommand{\LE}{\mathcal{L}(E)}
\newcommand{\rg}{\operatorname{rg}}
\newcommand{\vect}{\operatorname{Vect}}
\newcommand{\Mat}{\operatorname{Mat}}
\newcommand{\Id}{\operatorname{Id}}
\newcommand{\sh}{\operatorname{sh}}
\newcommand{\ch}{\operatorname{ch}}
\renewcommand{\th}{\operatorname{th}}
\DeclareMathOperator{\argsh}{argsh}
\newcommand{\vv}[1]{\overrightarrow{#1}}
\newcommand{\ve}[1]{\vec{#1}}
\newcommand{\ient}[2]{\llbracket #1,#2\rrbracket}
\newcommand{\dt}[1]{\left|\begin{matrix}#1\end{matrix}\right|}
\newcommand{\dis}{\displaystyle}

% ---------- Figures ----------
\pgfplotsset{trig format plots=rad}
\pgfplotsset{courbe/.style={axis lines=middle,axis equal image,
  axis line style={Gris,-{Stealth[length=5pt]}},
  xlabel={$x$},ylabel={$y$},ticklabel style={font=\scriptsize\sffamily,color=Gris},
  every axis x label/.style={at={(ticklabel* cs:1)},anchor=west,font=\small},
  every axis y label/.style={at={(ticklabel* cs:1)},anchor=south,font=\small},
  samples=300,smooth,width=0.6\linewidth}}
\tikzset{arcetude/.style={Bleu,line width=1.25pt},
         reste/.style={BleuClair,line width=1pt},
         tgte/.style={Ocre,line width=0.9pt},
         asy/.style={Gris,dashed,line width=0.6pt},
         pt/.style={circle,fill=Ink,inner sep=1.2pt}}

% ---------- Navigation ----------
\makeatletter
\newcommand*\l@rapentry{\@dottedtocline{1}{0em}{0em}}
\newcommand*\l@metentry{\@dottedtocline{1}{0em}{0em}}
\newcommand{\listemethodes}{\@starttoc{met}}
\newcommand{\listerappels}{\@starttoc{rap}}
\makeatother
% onglets de chapitre sur le bord droit
\usepackage{eso-pic}
\newif\ifonglets \ongletsfalse
\newcommand{\dessineronglet}{%
  \edef\ongcol{\ifnum\value{section}=3 Ocre\else Bleu\fi}%
  \edef\ongnum{\the\value{section}}%
  \pgfmathsetmacro{\ongy}{3.2+2.3*(\ongnum-1)}%
  \begin{tikzpicture}[remember picture,overlay]
  \fill[\ongcol] ([xshift=-0.75cm,yshift=-\ongy cm]current page.north east)
     rectangle ([yshift=-\ongy cm-1.9cm]current page.north east);
  \node[text=white,font=\sffamily\bfseries\large] at ([xshift=-0.37cm,yshift=-\ongy cm-0.95cm]current page.north east)
     {\ongnum};
  \end{tikzpicture}}
\AddToShipoutPictureBG{\ifonglets\ifnum\value{section}>0 \dessineronglet\fi\fi}
% blocs de navigation et d'enseignement
\definirbloc{commencer}{Gris}{Si vous commencez ici}
\definirbloc{horsprogramme}{Brique}{Hors programme}
\newcommand{\HP}{\,{\sffamily\bfseries\scriptsize\textcolor{Brique}{[HP]}}}
\newcommand{\source}[1]{{\footnotesize\sffamily\color{Gris}\,(#1)}}
\newcommand{\coupdepouce}[1]{\par\vspace{-0.2em}\noindent{\footnotesize\sffamily\color{Gris}%
  \textbf{Coup de pouce} (à donner si l'étudiant bloque) : #1}\par}
\newcommand{\pointarret}[1]{\par\vspace{1em}\noindent{\color{Filet}\rule[0.5ex]{\linewidth}{0.4pt}}\par\nobreak
  \noindent{\footnotesize\sffamily\color{Ocre}\bfseries POINT D'ARRÊT POSSIBLE}\quad{\footnotesize\sffamily\color{Gris}#1}\par\vspace{0.6em}}
\newcommand{\demarche}[1]{\noindent{\sffamily\bfseries\small\color{BleuMoyen}Démarche.}\ \textit{#1}\par\smallskip}

% ---------- Corrigés : idées, astuces, pièges, à retenir ----------
\newcommand{\choix}[1]{\par\noindent{\sffamily\bfseries\small\color{BleuMoyen}Choix de la méthode.}\ #1\par\smallskip}
\newcommand{\gain}[1]{\par\noindent{\sffamily\bfseries\small\color{Ocre}Astuce.}\ #1\par}
\newcommand{\piege}[1]{\par\noindent{\sffamily\bfseries\small\color{Brique}Piège.}\ #1\par}
\newtcolorbox{retenirbox}{enhanced,breakable,colback=Fond,colframe=Fond,arc=0pt,boxrule=0pt,
  borderline west={2pt}{0pt}{Ocre},left=8pt,right=6pt,top=4pt,bottom=4pt,before skip=0.7em,after skip=0.5em}
\newcommand{\retenir}[1]{\begin{retenirbox}{\sffamily\bfseries\footnotesize\color{Ocre}À RETENIR}\quad #1\end{retenirbox}}
\newcommand{\TD}[1]{{\color{Ocre}\bfseries TD\,#1}\ ·\ }
\newcommand{\q}[1]{\par\smallskip\noindent{\sffamily\bfseries\color{Bleu}#1}\ }
\providecommand{\iddots}{\mathinner{\mkern1mu\raise1pt\hbox{.}\mkern2mu\raise4pt\hbox{.}\mkern2mu\raise7pt\hbox{.}\mkern1mu}}


\begin{document}
\pagenumbering{roman}
\begin{titlepage}
\thispagestyle{empty}
\begin{tikzpicture}[remember picture,overlay]
  \fill[Bleu] (current page.north west) rectangle ([yshift=-15.5cm]current page.north east);
  \node[anchor=north west,text=white,font=\sffamily\small] at ([xshift=2.6cm,yshift=-3.2cm]current page.north west)
     {\textls[250]{MATHÉMATIQUES\ \ ·\ \ CPGE PT\ \ ·\ \ DEUXIÈME ANNÉE}};
  \node[anchor=north west,text=white] at ([xshift=2.5cm,yshift=-4.4cm]current page.north west)
     {\fontsize{62}{66}\selectfont\scshape Le Collé};
  \draw[Ocre,line width=1.6pt] ([xshift=2.6cm,yshift=-7.3cm]current page.north west) -- ++(4.2cm,0);
  \node[anchor=north west,text=white,font=\sffamily\Large\bfseries,align=left]
     at ([xshift=2.6cm,yshift=-8.0cm]current page.north west) {Algèbre linéaire et réduction};
  \node[anchor=north west,text=white!90!Bleu,font=\sffamily\large,align=left]
     at ([xshift=2.6cm,yshift=-9.0cm]current page.north west) {Sommes directes, sous-espaces stables, trace\\[0.15cm]Éléments propres, diagonalisation, trigonalisation};
  \node[anchor=north west,text=Ink,font=\sffamily\small,align=left,text width=12cm]
     at ([xshift=2.6cm,yshift=-17.4cm]current page.north west)
     {\textcolor{Bleu}{\bfseries Cours, méthodes et exercices corrigés}\\[0.3cm]\textcolor{Gris}{Les compléments d'algèbre linéaire de deuxième année, puis la réduction des endomorphismes et des matrices carrées.}};
  \node[anchor=south west,text=Ink,font=\sffamily\large] at ([xshift=2.6cm,yshift=3.0cm]current page.south west) {Ahmed Amine Ouqlani};
  \draw[Filet] ([xshift=2.6cm,yshift=2.6cm]current page.south west) -- ([xshift=-2.4cm,yshift=2.6cm]current page.south east);
\end{tikzpicture}
\end{titlepage}
\thispagestyle{plain}
\section*{Avant-propos}
\markboth{Avant-propos}{}
\providecommand{\Ker}{\operatorname{Ker}}\renewcommand{\Im}{\operatorname{Im}}\providecommand{\tr}{\operatorname{tr}}\providecommand{\Sp}{\operatorname{Sp}}
Ce livre accompagne les séances sur l'algèbre linéaire de deuxième année. Il reprend, dans l'ordre et avec les mêmes
notations, la présentation projetée en séance, et sert de référence pour travailler seul ensuite : chaque résultat y est
énoncé avec ses hypothèses, chaque méthode est illustrée par un exemple rédigé, et chaque partie se termine par quelques
exercices, corrigés au chapitre~2.

Le chapitre suit la progression du programme de PT. Les \emph{compléments d'algèbre linéaire} (matrices semblables, sommes
directes, trace, hyperplans, sous-espaces stables) fournissent le langage ; la \emph{réduction} l'utilise pour répondre à une
seule question : dans quelle base la matrice d'un endomorphisme est-elle la plus simple ? Éléments propres, polynôme
caractéristique, critères de diagonalisation, puis applications (puissances, suites, équations matricielles) et trigonalisation.
Le chapitre \emph{Déterminants} est supposé connu : il sert à calculer les polynômes caractéristiques.

Notations : $\chi_A(x)=\det(xI_n-A)$, $E_\lambda(f)=\Ker(f-\lambda\Id_E)$, $\Sp(f)$, $A=PDP^{-1}$ (colonnes de $P$ : vecteurs propres).
Les polynômes d'endomorphismes et les systèmes différentiels $X'=AX$ sont hors programme en PT ; les quelques compléments de ce
type sont signalés \HP{} ou placés dans un encadré « Pour aller plus loin ».
Les compléments (TD complet et ses corrections, sélection de cinq problèmes, fiche de référence) suivent la même numérotation.

\vfill
{\footnotesize\sffamily\color{Gris}\textcopyright~2026 Ahmed Amine Ouqlani. Tous droits réservés.\par}
\clearpage
\thispagestyle{plain}
\hypertarget{toc}{}
\renewcommand{\contentsname}{Sommaire}
\tableofcontents
\clearpage
\pagenumbering{arabic}
\ongletstrue
% =====================================================================
%       CHAPITRE 1 — ALGÈBRE LINÉAIRE : COMPLÉMENTS ET RÉDUCTION
% =====================================================================
\clearpage
\section{Algèbre linéaire et réduction}\label{sec:red}

Un endomorphisme se lit à travers sa matrice ; mais cette matrice dépend de la base choisie, et certaines bases rendent
tout simple. Prenons l'endomorphisme $f$ de $\R^3$ de matrice, dans la base canonique,
\[A=\begin{pmatrix}3&1&1\\2&3&2\\-2&-1&0\end{pmatrix}.\]
Les vecteurs $u_1=(0,-1,1)$, $u_2=(-1,0,1)$ et $u_3=(-1,-1,1)$ vérifient $f(u_1)=u_1$, $f(u_2)=2u_2$, $f(u_3)=3u_3$ (le vérifier) ;
ils forment une base (déterminant non nul), dans laquelle la matrice de $f$ est $D=\mathrm{diag}(1,2,3)$. Dans cette base,
$f$ agit sur chaque axe par une simple dilatation : $f^k$ a pour matrice $\mathrm{diag}(1,2^k,3^k)$, et $A^k=PD^kP^{-1}$ se
calcule sans effort.

\emph{Réduire} un endomorphisme, c'est chercher une telle base. Trois questions organisent le chapitre : quels vecteurs sont
simplement dilatés par $f$ (\emph{vecteurs propres}) ? Comment les trouver (\emph{polynôme caractéristique}) ? Quand en
trouve-t-on assez pour former une base (\emph{diagonalisation}) ? Les compléments d'algèbre linéaire du début (sommes directes,
sous-espaces stables, trace) fournissent le langage nécessaire.

\ouverture{%
\begin{itemize}[leftmargin=1.1em,itemsep=1pt]
\item Manipuler les sommes directes, les sous-espaces stables, la trace et les hyperplans.
\item Déterminer les valeurs propres et les sous-espaces propres d'un endomorphisme ou d'une matrice.
\item Décider si une matrice est diagonalisable, et la diagonaliser ($n=2$ ou $3$).
\item Utiliser la réduction : puissances, suites récurrentes, équations matricielles ; trace et déterminant en fonction des valeurs propres.
\end{itemize}}{%
Algèbre linéaire de première année : bases, dimension, théorème du rang, matrice d'une application linéaire, changement de base ;
chapitre \emph{Déterminants} ; polynômes (racines, multiplicité, factorisation).}{Trois séances}

Dans tout le chapitre, $\K=\R$ ou $\C$, $E$ est un $\K$-espace vectoriel et $f\in\LE$ un endomorphisme de $E$. Quand $E$ est de
dimension finie, on note $n=\dim E$. Pour une base $\mathcal B$ de $E$, $\Mat_{\mathcal B}(f)$ est la matrice de $f$ dans $\mathcal B$.

% ---------------------------------------------------------------------
\subsection{Compléments : matrices semblables, sommes directes, trace, hyperplans}\label{ss:complements}

\subsubsection*{Changement de base, matrices semblables}

\begin{proposition}[changement de base]\label{p:chgt}
Soient $\mathcal B$ et $\mathcal B'$ deux bases de $E$ et $P$ la matrice de passage de $\mathcal B$ à $\mathcal B'$ (ses colonnes
sont les coordonnées des vecteurs de $\mathcal B'$ dans $\mathcal B$). Si $A=\Mat_{\mathcal B}(f)$ et $A'=\Mat_{\mathcal B'}(f)$, alors
\[A'=P^{-1}AP.\]
\end{proposition}

\begin{definition}
Deux matrices $A,B\in\Mn{\K}$ sont \emph{semblables} s'il existe $P\in\GLn{\K}$ telle que $B=P^{-1}AP$. Deux matrices semblables
représentent le même endomorphisme dans deux bases.
\end{definition}

\begin{reflexe}
Tout ce qui ne dépend que de l'endomorphisme (et pas de la base) est commun à deux matrices semblables : rang, déterminant,
trace (ci-dessous), polynôme caractéristique et valeurs propres (§\ref{ss:polcar}). Inversement, si deux matrices diffèrent par
l'un de ces invariants, elles ne sont pas semblables.
\end{reflexe}

\subsubsection*{Somme de plusieurs sous-espaces, somme directe}

\begin{definition}
Soient $F_1,\dots,F_p$ des sous-espaces vectoriels de $E$. Leur \emph{somme} est
$F_1+\dots+F_p=\{x_1+\dots+x_p,\ x_i\in F_i\}$ : c'est un sous-espace vectoriel. La somme est \emph{directe}, et l'on note
$F_1\oplus\dots\oplus F_p$, si tout vecteur de la somme s'écrit de manière \emph{unique} $x_1+\dots+x_p$ avec $x_i\in F_i$.
\end{definition}

\begin{proposition}[caractérisation]\label{p:directe}
La somme $F_1+\dots+F_p$ est directe si et seulement si
\[x_1+\dots+x_p=0\text{ avec }x_i\in F_i\quad\Longrightarrow\quad x_1=\dots=x_p=0.\]
Pour deux sous-espaces : $F+G$ est directe si et seulement si $F\cap G=\{0\}$.
\end{proposition}

\begin{proof}
Si la somme est directe, $0=0+\dots+0$ est l'unique décomposition de $0$. Réciproquement, si $x_1+\dots+x_p=y_1+\dots+y_p$,
alors $(x_1-y_1)+\dots+(x_p-y_p)=0$ avec $x_i-y_i\in F_i$, donc $x_i=y_i$. Pour $p=2$ : si $u\in F\cap G$, $u+(-u)=0$ avec
$u\in F$, $-u\in G$.
\end{proof}

\begin{attention}
Pour $p\geq3$, il ne suffit pas que les intersections deux à deux soient nulles : dans $\R^2$, les trois droites
$\vect(1,0)$, $\vect(0,1)$, $\vect(1,1)$ se coupent deux à deux en $\{0\}$, mais $(1,0)+(0,1)-(1,1)=0$ : la somme n'est pas directe.
La seule caractérisation au programme est celle de la proposition~\ref{p:directe}.
\end{attention}

\begin{theoreme}[base adaptée, dimension]\label{t:adaptee}
On suppose $E$ de dimension finie. Si la somme $F_1+\dots+F_p$ est directe, la concaténation de bases des $F_i$ est une base
de $F_1\oplus\dots\oplus F_p$ (base \emph{adaptée}), et
\[\dim(F_1\oplus\dots\oplus F_p)=\dim F_1+\dots+\dim F_p.\]
En particulier, $E=F_1\oplus\dots\oplus F_p$ si et seulement si la concaténation de bases des $F_i$ est une base de $E$.
\end{theoreme}

\begin{proof}
La famille concaténée engendre la somme. Elle est libre : une combinaison linéaire nulle s'écrit $x_1+\dots+x_p=0$, où $x_i$
regroupe les termes de la base de $F_i$ ; la somme étant directe, chaque $x_i$ est nul, puis chaque coefficient (base de $F_i$).
\end{proof}

\begin{enpratique}
Pour montrer que $E=F\oplus G$ en dimension finie, il suffit de deux des trois propriétés : $F\cap G=\{0\}$, $F+G=E$,
$\dim F+\dim G=\dim E$. Le plus souvent : intersection nulle et somme des dimensions.
\end{enpratique}

\subsubsection*{Trace}

\begin{definition}
La \emph{trace} de $A=(a_{i,j})\in\Mn{\K}$ est la somme de ses coefficients diagonaux : $\tr(A)=\sum_{i=1}^na_{i,i}$.
\end{definition}

\begin{proposition}\label{p:trace}
\begin{enumerate}
\item La trace est une forme linéaire sur $\Mn{\K}$.
\item Pour $A\in\mathcal M_{n,p}(\K)$ et $B\in\mathcal M_{p,n}(\K)$ : $\tr(AB)=\tr(BA)$.
\item Deux matrices semblables ont la même trace. On définit donc la trace d'un endomorphisme $f$ en dimension finie :
$\tr(f)=\tr\big(\Mat_{\mathcal B}(f)\big)$, indépendante de la base $\mathcal B$.
\end{enumerate}
\end{proposition}

\begin{proof}
2. $\tr(AB)=\sum_{i=1}^n\sum_{k=1}^pa_{i,k}b_{k,i}=\sum_{k=1}^p\sum_{i=1}^nb_{k,i}a_{i,k}=\tr(BA)$.
3. $\tr(P^{-1}AP)=\tr\big((P^{-1}A)P\big)=\tr\big(P(P^{-1}A)\big)=\tr(A)$.
\end{proof}

\begin{attention}
$\tr(AB)\neq\tr(A)\tr(B)$ en général, et $\tr(ABC)=\tr(CAB)$ (permutation circulaire) mais $\tr(ABC)\neq\tr(BAC)$ en général.
\end{attention}

\begin{exemple}[trace d'un projecteur]
Si $p$ est un projecteur de rang $r$, dans une base adaptée à $E=\Im p\oplus\Ker p$, sa matrice est
$\mathrm{diag}(I_r,0)$ : $\tr(p)=r=\rg(p)$. Pour une symétrie $s$ : $\tr(s)=\dim\Ker(s-\Id)-\dim\Ker(s+\Id)$.
\end{exemple}

\subsubsection*{Hyperplans}

\begin{definition}
Un \emph{hyperplan} de $E$ ($\dim E=n\geq1$) est un sous-espace qui admet une droite comme supplémentaire ; c'est un
sous-espace de dimension $n-1$.
\end{definition}

\begin{proposition}\label{p:hyperplan}
\begin{enumerate}
\item Les hyperplans de $E$ sont les noyaux des formes linéaires non nulles. Dans une base, un hyperplan a une équation
$a_1x_1+\dots+a_nx_n=0$ avec $(a_1,\dots,a_n)\neq(0,\dots,0)$, unique à un facteur non nul près.
\item Si $H=\Ker\varphi$ et $u\notin H$, alors $E=H\oplus\vect(u)$.
\item L'intersection de $p$ hyperplans est de dimension au moins $n-p$. Réciproquement, tout sous-espace de dimension $n-p$
est l'intersection de $p$ hyperplans : il est défini par un système de $p$ équations.
\end{enumerate}
\end{proposition}

\begin{enpratique}
L'ensemble des solutions d'un système linéaire homogène de $p$ équations à $n$ inconnues est un sous-espace de dimension
$n-r$, où $r$ est le rang du système (théorème du rang) : chaque équation « indépendante » fait perdre une dimension.
Dans $\R^3$, un plan vectoriel a une équation, une droite vectorielle deux.
\end{enpratique}

\begin{exo}{red-somme}{1}{sommes directes}
\begin{enumerate}
\item Dans $\R^3$, on pose $F=\{(x,y,z),\ x+y+z=0\}$ et $G=\vect\big((1,1,1)\big)$. Montrer que $\R^3=F\oplus G$ et décomposer
$v=(x,y,z)$ selon cette somme.
\item Dans $\R^3$, les droites $D_1=\vect(1,0,0)$, $D_2=\vect(0,1,0)$, $D_3=\vect(1,1,0)$ sont-elles en somme directe ?
\end{enumerate}
\end{exo}
\begin{exo}{red-trace}{2}{trace}
\begin{enumerate}
\item Montrer qu'il n'existe pas de matrices $A,B\in\Mn{\K}$ telles que $AB-BA=I_n$.
\item Soit $A\in\Mn{\K}$ telle que $A^2=A$. Montrer que $\rg(A)=\tr(A)$.
\end{enumerate}
\end{exo}
\begin{exo}{red-hyp}{2}{hyperplans de $\R^4$}
Dans $\R^4$, on considère $H_1$ : $x+y+z+t=0$ et $H_2$ : $x-y+z-t=0$, et $u=(1,1,1,1)$.
\begin{enumerate}
\item Donner la dimension et une base de $H_1\cap H_2$.
\item La droite $\vect(u)$ est-elle un supplémentaire de $H_1$ ? de $H_2$ ?
\end{enumerate}
\end{exo}

% ---------------------------------------------------------------------
\subsection{Sous-espaces stables}\label{ss:stables}

\begin{definition}
Un sous-espace $F$ de $E$ est \emph{stable} par $f$ si $f(F)\subset F$. L'application $f_F:F\to F$, $x\mapsto f(x)$, est alors un
endomorphisme de $F$, l'\emph{endomorphisme induit} par $f$ sur $F$.
\end{definition}

\begin{proposition}\label{p:stable}
\begin{enumerate}
\item Si $F=\vect(e_1,\dots,e_p)$, $F$ est stable par $f$ si et seulement si $f(e_j)\in F$ pour tout $j$ : on vérifie sur une
famille génératrice.
\item $\Ker f$ et $\Im f$ sont stables par $f$. Plus généralement, si $f\circ g=g\circ f$, alors $\Ker g$ et $\Im g$ sont stables par $f$.
\end{enumerate}
\end{proposition}

\begin{proof}
2. Si $g(x)=0$, $g\big(f(x)\big)=f\big(g(x)\big)=0$ : $f(x)\in\Ker g$. Si $y=g(x)$, $f(y)=g\big(f(x)\big)\in\Im g$. Pour $g=f$, on
obtient la stabilité de $\Ker f$ et $\Im f$.
\end{proof}

\begin{theoreme}[traduction matricielle]\label{t:blocs}
On suppose $E$ de dimension finie. Soit $F$ un sous-espace de dimension $p$ et $\mathcal B=(e_1,\dots,e_n)$ une base adaptée à $F$
($F=\vect(e_1,\dots,e_p)$).
\begin{enumerate}
\item $F$ est stable par $f$ si et seulement si $\Mat_{\mathcal B}(f)$ est triangulaire par blocs :
$\begin{pmatrix}A&B\\0&C\end{pmatrix}$, avec $A\in\mathcal M_p(\K)$ ; $A$ est alors la matrice de l'endomorphisme induit $f_F$.
\item Si $E=F\oplus G$ et si $\mathcal B$ est adaptée à cette somme, $F$ et $G$ sont stables par $f$ si et seulement si
$\Mat_{\mathcal B}(f)$ est diagonale par blocs : $\begin{pmatrix}A&0\\0&C\end{pmatrix}$.
\end{enumerate}
Dans les deux cas, $\det f=\det A\times\det C$ et $\tr f=\tr A+\tr C$.
\end{theoreme}

\begin{proof}
La $j$-ième colonne contient les coordonnées de $f(e_j)$. Pour $j\leq p$, $f(e_j)\in F$ si et seulement si ses coordonnées sur
$e_{p+1},\dots,e_n$ sont nulles : c'est le bloc $0$. Le déterminant d'une matrice triangulaire par blocs est le produit des
déterminants des blocs diagonaux (développer selon les premières colonnes, chapitre \emph{Déterminants}).
\end{proof}

\begin{enpratique}
Lire une matrice par blocs, c'est lire des sous-espaces stables, et réciproquement. Un zéro en bas à gauche dit : « les $p$
premiers vecteurs de base engendrent un sous-espace stable ». Un projecteur et une symétrie sont les premiers exemples :
\[p\text{ projecteur : }E=\Im p\oplus\Ker p,\quad\Mat(p)=\begin{pmatrix}I_r&0\\0&0\end{pmatrix};\]
\[s\text{ symétrie : }E=\Ker(s-\Id)\oplus\Ker(s+\Id),\quad\Mat(s)=\begin{pmatrix}I_r&0\\0&-I_{n-r}\end{pmatrix}.\]
\end{enpratique}

\begin{proposition}[droites stables]\label{p:droite}
Soit $u\neq0$. La droite $\vect(u)$ est stable par $f$ si et seulement s'il existe $\lambda\in\K$ tel que $f(u)=\lambda u$.
\end{proposition}

\begin{proof}
$\vect(u)$ est stable si et seulement si $f(u)\in\vect(u)$, c'est-à-dire $f(u)=\lambda u$ pour un scalaire $\lambda$.
\end{proof}

Cette dernière proposition fait la transition : chercher les droites stables, c'est chercher les vecteurs que $f$ se contente de
dilater. S'il y en a assez pour former une base, la matrice de $f$ dans cette base est diagonale.

\begin{exo}{red-stable}{2}{lire des blocs}
Soit $f$ l'endomorphisme de $\R^3$ de matrice $A=\begin{pmatrix}3&-1&1\\0&2&0\\1&-1&3\end{pmatrix}$ dans la base canonique
$(e_1,e_2,e_3)$.
\begin{enumerate}
\item Montrer que le plan $F$ d'équation $y=0$ et la droite $G=\vect(u)$, $u=(1,1,0)$, sont stables par $f$.
\item Montrer que $\R^3=F\oplus G$ et écrire la matrice de $f$ dans la base $(e_1,e_3,u)$. En déduire $\det A$ et $\tr A$.
\end{enumerate}
\end{exo}

% ---------------------------------------------------------------------
\subsection{Éléments propres}\label{ss:propres}

\begin{definition}
Soit $f\in\LE$ ($E$ de dimension quelconque).
\begin{itemize}
\item $\lambda\in\K$ est une \emph{valeur propre} de $f$ s'il existe $x\neq0$ tel que $f(x)=\lambda x$ ;
\item un tel $x$, \emph{non nul}, est un \emph{vecteur propre} associé à $\lambda$ ;
\item le \emph{sous-espace propre} associé à $\lambda$ est $E_\lambda(f)=\Ker(f-\lambda\Id_E)$.
\end{itemize}
En dimension finie, l'ensemble des valeurs propres est le \emph{spectre} de $f$, noté $\Sp(f)$.
\end{definition}

\begin{reflexe}
$E_\lambda(f)$ est toujours un sous-espace vectoriel ; il est différent de $\{0\}$ exactement quand $\lambda$ est valeur propre.
Il contient $0$ et les vecteurs propres associés à $\lambda$. En particulier, $0$ est valeur propre si et seulement si $f$ n'est pas
injective, et $E_0(f)=\Ker f$. Chaque sous-espace propre est stable par $f$ (sur $E_\lambda$, $f$ agit comme $\lambda\Id$).
\end{reflexe}

\begin{exemple}[la dérivation]\label{ex:derivation}
Sur $E=\mathcal C^\infty(\R,\R)$, soit $D:y\mapsto y'$. L'équation $D(y)=\lambda y$ est l'équation différentielle $y'=\lambda y$,
dont les solutions sont $t\mapsto Ke^{\lambda t}$ (chapitre \emph{Équations différentielles}). Tout réel $\lambda$ est valeur propre,
et $E_\lambda(D)=\vect(t\mapsto e^{\lambda t})$ est une droite. En dimension infinie, il peut y avoir une infinité de valeurs propres.
\end{exemple}

\begin{exemple}[un endomorphisme de polynômes]
Sur $\R_n[X]$, soit $f(P)=XP'$. Alors $f(X^k)=kX^k$ : chaque $X^k$ est vecteur propre, associé à la valeur propre $k$. La base
canonique est formée de vecteurs propres, la matrice de $f$ y est $\mathrm{diag}(0,1,\dots,n)$, et $\Sp(f)=\{0,1,\dots,n\}$.
\end{exemple}

\begin{theoreme}[sous-espaces propres en somme directe]\label{t:somme-propres}
Si $\lambda_1,\dots,\lambda_p$ sont des valeurs propres \emph{distinctes} de $f$, la somme $E_{\lambda_1}(f)+\dots+E_{\lambda_p}(f)$
est directe. De manière équivalente : toute famille de vecteurs propres associés à des valeurs propres distinctes est libre.
\end{theoreme}

\begin{proof}
Par récurrence sur $p$ ; $p=1$ est clair. Supposons le résultat vrai pour $p-1$ et soit $x_1+\dots+x_p=0$ avec $x_i\in E_{\lambda_i}$.
On applique $f$ : $\lambda_1x_1+\dots+\lambda_px_p=0$ ; on retranche $\lambda_p$ fois la première relation :
\[(\lambda_1-\lambda_p)x_1+\dots+(\lambda_{p-1}-\lambda_p)x_{p-1}=0.\]
Par hypothèse de récurrence, $(\lambda_i-\lambda_p)x_i=0$ pour $i<p$, donc $x_i=0$ car $\lambda_i\neq\lambda_p$ ; puis $x_p=0$.
\end{proof}

\begin{corollaire}\label{c:n-valeurs}
En dimension $n$, un endomorphisme a au plus $n$ valeurs propres.
\end{corollaire}

\begin{definition}
Pour $A\in\Mn{\K}$, $\lambda\in\K$ est valeur propre de $A$ s'il existe une colonne $X\neq0$ telle que $AX=\lambda X$ ;
$E_\lambda(A)=\{X\in\mathcal M_{n,1}(\K),\ AX=\lambda X\}$ ; $\Sp(A)$ est l'ensemble des valeurs propres.
Les éléments propres de $A$ sont ceux de l'endomorphisme de $\K^n$ canoniquement associé ; si $A=\Mat_{\mathcal B}(f)$, ce sont
ceux de $f$, lus en coordonnées dans $\mathcal B$.
\end{definition}

\begin{methode}{Déterminer un sous-espace propre}\label{m:sep}
Pour $\lambda$ donné, résoudre le système homogène $(A-\lambda I_n)X=0$. Par le théorème du rang,
\[\dim E_\lambda(A)=n-\rg(A-\lambda I_n).\]
On commence souvent par le rang : il donne la dimension avant tout calcul. $\lambda$ est valeur propre si et seulement si
$A-\lambda I_n$ n'est pas inversible, c'est-à-dire $\det(A-\lambda I_n)=0$.
\end{methode}

\begin{attention}
Une matrice réelle peut n'avoir aucune valeur propre réelle : la rotation $R=\begin{pmatrix}0&-1\\1&0\end{pmatrix}$ ne fixe aucune
droite du plan. Vue dans $\mathcal M_2(\C)$, elle a les valeurs propres $i$ et $-i$. On précise donc toujours le corps :
$\Sp_\R(A)\subset\Sp_\C(A)$.
\end{attention}

\begin{exo}{red-elem}{2}{éléments propres d'un endomorphisme de polynômes}
Soit $n\geq1$ et $f:\R_n[X]\to\R_n[X]$, $P\mapsto(X+1)P'$.
\begin{enumerate}
\item Vérifier que $f$ est un endomorphisme de $\R_n[X]$.
\item Montrer que $(X+1)^k$ est un vecteur propre de $f$ pour $0\leq k\leq n$. En déduire $\Sp(f)$ et les sous-espaces propres.
\end{enumerate}
\end{exo}

% ---------------------------------------------------------------------
\subsection{Polynôme caractéristique}\label{ss:polcar}

\begin{proposition}[polynôme caractéristique]\label{p:polcar}
Pour $A\in\Mn{\K}$, la fonction $x\mapsto\det(xI_n-A)$ est polynomiale, unitaire, de degré $n$. On note
\[\chi_A(x)=\det(xI_n-A)\]
le \emph{polynôme caractéristique} de $A$. On a de plus $\chi_A(x)=x^n-\tr(A)\,x^{n-1}+\dots+(-1)^n\det A$ ; pour $n=2$ :
\[\chi_A(x)=x^2-\tr(A)\,x+\det A.\]
\end{proposition}

{\small La démonstration (développement selon une colonne, par récurrence) n'est pas exigible pour $n\geq4$. Le coefficient
constant est $\chi_A(0)=\det(-A)=(-1)^n\det A$ ; celui de $x^{n-1}$ sera justifié au §\ref{ss:trigo}.}

\begin{theoreme}\label{t:racines}
\begin{enumerate}
\item Les valeurs propres de $A$ sont les racines de $\chi_A$ dans $\K$.
\item Deux matrices semblables ont le même polynôme caractéristique. On définit $\chi_f=\chi_{\Mat_{\mathcal B}(f)}$, indépendant
de la base, et $\Sp(f)$ est l'ensemble des racines de $\chi_f$ dans $\K$.
\item Si $A$ est triangulaire, $\chi_A(x)=\prod_{i=1}^n(x-a_{i,i})$ : ses valeurs propres sont ses coefficients diagonaux.
\end{enumerate}
\end{theoreme}

\begin{proof}
1. $\lambda\in\Sp(A)\iff\lambda I_n-A$ non inversible $\iff\det(\lambda I_n-A)=0$.
2. $xI_n-P^{-1}AP=P^{-1}(xI_n-A)P$, de même déterminant que $xI_n-A$.
3. $xI_n-A$ est triangulaire, de diagonale $(x-a_{i,i})$.
\end{proof}

\begin{definition}
La \emph{multiplicité} $m(\lambda)$ d'une valeur propre $\lambda$ est sa multiplicité comme racine de $\chi_f$. On dit que $\chi_f$
est \emph{scindé} sur $\K$ s'il est produit de facteurs de degré $1$ : $\chi_f(x)=\prod_{\lambda\in\Sp(f)}(x-\lambda)^{m(\lambda)}$,
et alors $\sum m(\lambda)=n$. Sur $\C$, tout polynôme est scindé (théorème de d'Alembert-Gauss).
\end{definition}

\begin{theoreme}[dimension et multiplicité]\label{t:dim-mult}
Pour toute valeur propre $\lambda$ de $f$ : \quad $1\leq\dim E_\lambda(f)\leq m(\lambda)$.
En particulier, si $\lambda$ est une valeur propre simple, $\dim E_\lambda(f)=1$.
\end{theoreme}

\begin{proof}
Soit $d=\dim E_\lambda(f)\geq1$. Dans une base adaptée à $E_\lambda(f)$, qui est stable, la matrice de $f$ est
$\begin{pmatrix}\lambda I_d&B\\0&C\end{pmatrix}$ (théorème~\ref{t:blocs}). Alors $\chi_f(x)=(x-\lambda)^d\chi_C(x)$ : $(x-\lambda)^d$
divise $\chi_f$, donc $d\leq m(\lambda)$.
\end{proof}

\begin{methode}{Calculer un polynôme caractéristique}\label{m:polcar}
\begin{enumerate}
\item $n=2$ : $\chi_A=x^2-\tr(A)\,x+\det A$, sans calcul de déterminant.
\item $n=3$ : calculer $\det(xI_3-A)$ en cherchant d'abord une structure (chapitre \emph{Déterminants}) : somme des lignes
constante, opération qui fait apparaître un facteur $(x-\lambda)$, développement selon une ligne contenant des zéros.
\emph{Factoriser au fur et à mesure} : on veut les racines, pas un polynôme développé.
\item Contrôles : $\chi_A(0)=(-1)^n\det A$ ; somme des racines (avec multiplicité) $=\tr A$ ; produit $=\det A$ (§\ref{ss:trigo}).
\end{enumerate}
\end{methode}

\begin{exemple}
$J=\begin{pmatrix}0&1&1\\1&0&1\\1&1&0\end{pmatrix}$. Dans $xI_3-J$, chaque ligne a pour somme $x-2$ : avec $C_1\leftarrow C_1+C_2+C_3$
puis $L_2\leftarrow L_2-L_1$, $L_3\leftarrow L_3-L_1$,
\[\chi_J(x)=(x-2)\dt{1&-1&-1\\1&x&-1\\1&-1&x}=(x-2)\dt{1&-1&-1\\0&x+1&0\\0&0&x+1}=(x-2)(x+1)^2.\]
Contrôle : $2-1-1=0=\tr J$ et $2\times(-1)^2=2=\det J$. Valeurs propres : $2$ (simple) et $-1$ (double).
\end{exemple}

\begin{attention}
On définit $\chi_A(x)=\det(xI_n-A)$, unitaire. Le calcul de $\det(A-xI_n)=(-1)^n\chi_A(x)$ donne les mêmes racines, mais n'est pas
le polynôme caractéristique quand $n$ est impair.
\end{attention}

\begin{exo}{red-chi}{2}{matrice compagnon}
Soient $a,b,c\in\K$ et $C=\begin{pmatrix}0&0&-c\\1&0&-b\\0&1&-a\end{pmatrix}$. Montrer que $\chi_C(x)=x^3+ax^2+bx+c$.
En déduire une matrice dont les valeurs propres sont $1$, $-1$ et $2$.
\end{exo}
% ---------------------------------------------------------------------
\subsection{Diagonalisation}\label{ss:diag}

\begin{definition}
On suppose $E$ de dimension finie.
\begin{itemize}
\item $f$ est \emph{diagonalisable} s'il existe une base de $E$ dans laquelle sa matrice est diagonale, c'est-à-dire une base
de $E$ formée de vecteurs propres de $f$.
\item $A\in\Mn{\K}$ est \emph{diagonalisable} si elle est semblable à une matrice diagonale : $A=PDP^{-1}$, avec $P$ inversible et
$D$ diagonale. Les colonnes de $P$ sont alors des vecteurs propres de $A$, et les coefficients de $D$ les valeurs propres
associées, dans le même ordre.
\end{itemize}
\end{definition}

\begin{theoreme}[critères de diagonalisabilité]\label{t:criteres}
Soit $f\in\LE$, $\dim E=n$. Les propriétés suivantes sont équivalentes :
\begin{enumerate}
\item $f$ est diagonalisable ;
\item $E=\bigoplus_{\lambda\in\Sp(f)}E_\lambda(f)$ ;
\item $\sum_{\lambda\in\Sp(f)}\dim E_\lambda(f)=n$ ;
\item $\chi_f$ est scindé sur $\K$ et, pour toute valeur propre $\lambda$, $\dim E_\lambda(f)=m(\lambda)$.
\end{enumerate}
\end{theoreme}

\begin{proof}
La somme des sous-espaces propres est directe (théorème~\ref{t:somme-propres}) ; sa dimension est donc $\sum\dim E_\lambda$
(théorème~\ref{t:adaptee}). (1)$\Rightarrow$(2) : une base de vecteurs propres est contenue dans la somme, qui vaut donc $E$.
(2)$\Leftrightarrow$(3) : un sous-espace de dimension $n$ est $E$. (2)$\Rightarrow$(1) : une base adaptée à la somme directe est
formée de vecteurs propres. (3)$\Leftrightarrow$(4) : on a toujours $\sum\dim E_\lambda\leq\sum m(\lambda)\leq\deg\chi_f=n$
(théorème~\ref{t:dim-mult}) ; il y a égalité partout si et seulement si $\sum m(\lambda)=n$ ($\chi_f$ scindé) et
$\dim E_\lambda=m(\lambda)$ pour tout $\lambda$.
\end{proof}

\begin{corollaire}\label{c:distinctes}
Si $f$ a $n$ valeurs propres distinctes ($\chi_f$ scindé à racines simples), $f$ est diagonalisable, et chaque sous-espace propre
est une droite. La réciproque est fausse : $\Id_E$ est diagonalisable avec une seule valeur propre.
\end{corollaire}

\begin{reflexe}
Une matrice diagonalisable qui n'a qu'une valeur propre $\lambda$ est égale à $\lambda I_n$ : $A=P(\lambda I_n)P^{-1}=\lambda I_n$.
Ainsi $\begin{pmatrix}1&1\\0&1\end{pmatrix}$ n'est pas diagonalisable : c'est le réflexe le plus rapide pour une matrice
triangulaire à diagonale constante.
\end{reflexe}

\begin{exemple}[projecteurs et symétries]
Un projecteur est diagonalisable : $E=\Im p\oplus\Ker p=E_1(p)\oplus E_0(p)$. Une symétrie aussi :
$E=E_1(s)\oplus E_{-1}(s)$. Ce sont les deux exemples géométriques du programme.
\end{exemple}

\begin{methode}{Diagonaliser une matrice ($n=2$ ou $3$)}\label{m:diag}
\begin{enumerate}
\item Calculer $\chi_A$ et le \emph{factoriser} (méthode~\ref{m:polcar}). S'il n'est pas scindé sur $\K$ : pas diagonalisable sur $\K$.
\item Si les racines sont simples : diagonalisable (corollaire~\ref{c:distinctes}) ; sinon, pour chaque racine multiple, comparer
$\dim E_\lambda=n-\rg(A-\lambda I_n)$ à $m(\lambda)$.
\item Déterminer une base de chaque $E_\lambda$ ; les concaténer en colonnes : c'est $P$. Placer les valeurs propres dans $D$ dans
le même ordre.
\item Conclure $A=PDP^{-1}$. Contrôle : vérifier $AP=PD$ (colonne par colonne : $AX=\lambda X$), sans calculer $P^{-1}$.
\end{enumerate}
\end{methode}

\begin{exemple}[trois valeurs propres distinctes]\label{ex:diag1}
$A=\begin{pmatrix}5&1&-1\\2&4&-2\\1&-1&3\end{pmatrix}$. Dans $xI_3-A$, on fait $C_2\leftarrow C_2+C_3$ : la deuxième colonne devient
$(0,x-2,x-2)$ ; on sort $x-2$, puis $L_3\leftarrow L_3-L_2$ et l'on développe selon la deuxième colonne :
\[\chi_A(x)=(x-2)\dt{x-5&0&1\\-2&1&2\\1&0&x-5}=(x-2)\big((x-5)^2-1\big)=(x-2)(x-4)(x-6).\]
Trois valeurs propres distinctes : $A$ est diagonalisable. En résolvant $(A-\lambda I_3)X=0$ :
$E_2=\vect(0,1,1)$, $E_4=\vect(1,0,1)$, $E_6=\vect(1,1,0)$. Donc
\[A=PDP^{-1},\qquad P=\begin{pmatrix}0&1&1\\1&0&1\\1&1&0\end{pmatrix},\quad D=\mathrm{diag}(2,4,6).\]
\end{exemple}

\begin{exemple}[une valeur propre double]\label{ex:diag2}
$B=\begin{pmatrix}3&-1&1\\0&2&0\\1&-1&3\end{pmatrix}$. En développant selon la deuxième ligne,
$\chi_B(x)=(x-2)\big((x-3)^2-1\big)=(x-2)^2(x-4)$. Pour la valeur double : $B-2I_3=\begin{pmatrix}1&-1&1\\0&0&0\\1&-1&1\end{pmatrix}$
est de rang $1$, donc $\dim E_2=2=m(2)$ : $B$ est diagonalisable. $E_2$ est le plan $x-y+z=0$, de base $(1,1,0)$, $(-1,0,1)$ ;
$E_4=\vect(1,0,1)$.
\end{exemple}

\begin{exemple}[deux matrices non diagonalisables]
$C=\begin{pmatrix}2&1\\-1&4\end{pmatrix}$ : $\chi_C=x^2-6x+9=(x-3)^2$ et $C-3I_2=\begin{pmatrix}-1&1\\-1&1\end{pmatrix}$ est de rang $1$ :
$\dim E_3=1<2$. La rotation $R=\begin{pmatrix}0&-1\\1&0\end{pmatrix}$ : $\chi_R=x^2+1$ n'est pas scindé sur $\R$ ; $R$ n'est pas
diagonalisable dans $\mathcal M_2(\R)$, mais elle l'est dans $\mathcal M_2(\C)$ (deux valeurs propres distinctes $\pm i$).
\end{exemple}

\begin{exo}{red-diag1}{1}{diagonaliser}
Diagonaliser $A=\begin{pmatrix}1&-1&-1\\-1&1&-1\\-1&-1&1\end{pmatrix}$ : donner $P$ inversible et $D$ diagonale telles que $A=PDP^{-1}$.
\end{exo}
\begin{exo}{red-diag2}{2}{diagonalisabilité selon un paramètre}
Pour $a\in\R$, on pose $A_a=\begin{pmatrix}1&a&1\\0&1&1\\0&0&2\end{pmatrix}$. Pour quelles valeurs de $a$ la matrice $A_a$ est-elle
diagonalisable ? La diagonaliser dans ce cas.
\end{exo}
\begin{exo}{red-rang1}{3}{la matrice dont tous les coefficients valent 1}
Soit $n\geq2$ et $J_n\in\Mn{\R}$ la matrice dont tous les coefficients valent $1$.
\begin{enumerate}
\item Déterminer le rang de $J_n$, puis ses valeurs propres et ses sous-espaces propres. Est-elle diagonalisable ?
\item Soient $a,b$ réels, $b\neq0$. En déduire les valeurs propres de $aI_n+bJ_n$, et son déterminant.
\end{enumerate}
\end{exo}

% ---------------------------------------------------------------------
\subsection{Applications de la diagonalisation}\label{ss:applications}

\begin{proposition}[puissances]\label{p:puissances}
Si $A=PDP^{-1}$ avec $D=\mathrm{diag}(\lambda_1,\dots,\lambda_n)$, alors pour tout $k\in\N$ :
\[A^k=PD^kP^{-1},\qquad D^k=\mathrm{diag}(\lambda_1^k,\dots,\lambda_n^k).\]
Si de plus $A$ est inversible (aucun $\lambda_i$ nul), la formule vaut pour $k\in\Z$ ; en particulier $A^{-1}=PD^{-1}P^{-1}$.
\end{proposition}

\begin{proof}
$A^k=(PDP^{-1})(PDP^{-1})\cdots(PDP^{-1})=PD^kP^{-1}$ : les produits $P^{-1}P$ intermédiaires se simplifient (récurrence).
\end{proof}

\begin{methode}{Suites récurrentes couplées}\label{m:suites}
Un système $u_{k+1}=au_k+bv_k$, $v_{k+1}=cu_k+dv_k$ s'écrit $X_{k+1}=AX_k$ avec $X_k=\begin{pmatrix}u_k\\v_k\end{pmatrix}$, donc
$X_k=A^kX_0$. Si $A=PDP^{-1}$, on pose $Y_k=P^{-1}X_k$ (coordonnées dans la base de vecteurs propres) : $Y_{k+1}=DY_k$, et
\emph{chaque coordonnée est une suite géométrique}. On revient à $X_k=PY_k$.
Une récurrence linéaire d'ordre $3$, $u_{k+3}=\alpha u_{k+2}+\beta u_{k+1}+\gamma u_k$, se traite de même avec
$X_k=(u_k,u_{k+1},u_{k+2})$ et une matrice compagnon (exercice~\ref{e:red-chi}).
\end{methode}

\begin{exemple}
$u_{k+1}=u_k+2v_k$, $v_{k+1}=2u_k+v_k$, $u_0=1$, $v_0=0$. Ici $A=\begin{pmatrix}1&2\\2&1\end{pmatrix}$, $\chi_A=x^2-2x-3=(x-3)(x+1)$,
$E_3=\vect(1,1)$, $E_{-1}=\vect(1,-1)$. Les coordonnées dans la base propre sont, à un facteur près, $s_k=u_k+v_k$ et
$d_k=u_k-v_k$ : on vérifie directement $s_{k+1}=3s_k$ et $d_{k+1}=-d_k$. Donc $s_k=3^k$, $d_k=(-1)^k$, et
\[u_k=\frac{3^k+(-1)^k}2,\qquad v_k=\frac{3^k-(-1)^k}2.\]
\end{exemple}

\begin{proposition}[commuter avec une matrice diagonale]\label{p:commute}
Si $D=\mathrm{diag}(d_1,\dots,d_n)$ a des coefficients \emph{deux à deux distincts}, les matrices qui commutent avec $D$ sont les
matrices diagonales.
\end{proposition}

\begin{proof}
Le coefficient $(i,j)$ de $MD-DM$ vaut $m_{i,j}(d_j-d_i)$. Il est nul pour $i\neq j$ si et seulement si $m_{i,j}=0$.
\end{proof}

\begin{enpratique}
Pour résoudre une équation matricielle ($M^2=A$, $MA=AM$\dots) avec $A=PDP^{-1}$, on pose $N=P^{-1}MP$ : l'équation devient la
même équation avec $D$, beaucoup plus simple (exercice~\ref{e:red-racine}).
\end{enpratique}

\begin{plusloin}
Le système différentiel $X'=AX$ (hors programme en PT) se traite comme les suites : avec $A=PDP^{-1}$ et $Y=P^{-1}X$, on obtient
$Y'=DY$, c'est-à-dire $n$ équations $y_i'=\lambda_iy_i$ indépendantes, résolues au chapitre \emph{Équations différentielles}.
\end{plusloin}

\begin{exo}{red-puiss}{2}{puissances et suites}
Soit $A=\begin{pmatrix}1&2\\2&1\end{pmatrix}$. Calculer $A^k$ pour $k\in\N$, puis exprimer $u_k$ et $v_k$ pour les suites définies par
$u_0=1$, $v_0=0$, $u_{k+1}=u_k+2v_k$, $v_{k+1}=2u_k+v_k$, en utilisant $A^k$. Retrouver le résultat de l'exemple du cours (§1.6).
\end{exo}
\begin{exo}{red-racine}{3}{racines carrées d'une matrice}
\begin{enumerate}
\item Déterminer les matrices $M\in\mathcal M_3(\R)$ telles que $M^2=D$, où $D=\mathrm{diag}(1,4,9)$. \emph{(Montrer d'abord que $M$ commute avec $D$.)}
\item Soit $A\in\mathcal M_3(\R)$ diagonalisable, de valeurs propres $1$, $4$, $9$. Combien existe-t-il de matrices $M$ telles que $M^2=A$ ?
\end{enumerate}
\end{exo}

% ---------------------------------------------------------------------
\subsection{Trigonalisation}\label{ss:trigo}

\begin{definition}
$f$ est \emph{trigonalisable} s'il existe une base dans laquelle sa matrice est triangulaire (supérieure). Une matrice est
trigonalisable si elle est semblable à une matrice triangulaire.
\end{definition}

\begin{theoreme}[admis]\label{t:trigo}
$f$ (ou $A$) est trigonalisable si et seulement si son polynôme caractéristique est scindé. En particulier, toute matrice de
$\Mn{\C}$ est trigonalisable.
\end{theoreme}

{\small La démonstration est hors programme, et aucune technique générale de trigonalisation n'est exigible : les exercices la
guident (exercice~\ref{e:red-trig}).}

\begin{proposition}[trace, déterminant et valeurs propres]\label{p:trdet}
Si $\chi_A$ est scindé, $\chi_A(x)=\prod_{i=1}^n(x-\lambda_i)$ (valeurs propres répétées selon leur multiplicité), alors
\[\tr(A)=\lambda_1+\dots+\lambda_n,\qquad\det(A)=\lambda_1\cdots\lambda_n.\]
C'est toujours le cas dans $\Mn{\C}$.
\end{proposition}

\begin{proof}
$A$ est semblable à une matrice triangulaire $T$, qui a le même polynôme caractéristique : la diagonale de $T$ porte donc les
$\lambda_i$, avec leur multiplicité (théorème~\ref{t:racines}). Trace et déterminant sont invariants par similitude, et ceux de
$T$ sont la somme et le produit de ses coefficients diagonaux.
\end{proof}

\begin{astuce}
Trace et déterminant donnent la somme et le produit des valeurs propres \emph{sans calcul} : pour une matrice $2\times2$, ils
déterminent les deux valeurs propres ; pour une matrice $3\times3$ dont on connaît déjà deux valeurs propres, la trace donne la
troisième. Et si $\rg(A)=1$, $0$ est valeur propre d'ordre au moins $n-1$, et la dernière valeur propre est $\tr(A)$.
\end{astuce}

\begin{exemple}[une matrice nilpotente]
Si $A\in\Mn{\C}$ vérifie $A^p=0$ pour un entier $p\geq1$, et si $AX=\lambda X$ avec $X\neq0$, alors $0=A^pX=\lambda^pX$, donc
$\lambda=0$ : $\Sp(A)=\{0\}$, $\chi_A=x^n$, et $\tr A=0$. Une telle matrice n'est diagonalisable que si elle est nulle (réflexe du
§\ref{ss:diag}).
\end{exemple}

\begin{exo}{red-nilp}{2}{matrices nilpotentes}
Soit $A\in\Mn{\C}$ telle que $A^p=0$ pour un entier $p\geq1$.
\begin{enumerate}
\item Montrer que $\Sp(A)=\{0\}$, puis que $\chi_A(x)=x^n$. En déduire $\tr A$ et $\det A$.
\item Montrer que $A$ est diagonalisable si et seulement si $A=0$. Donner une matrice $A\neq0$ de $\mathcal M_2(\R)$ telle que $A^2=0$.
\end{enumerate}
\end{exo}
\begin{exo}{red-trig}{3}{trigonaliser une matrice guidée}
Soit $T=\begin{pmatrix}5&-17&25\\2&-9&16\\1&-5&9\end{pmatrix}$.
\begin{enumerate}
\item Calculer $\chi_T$ et montrer que $T$ n'est pas diagonalisable, même dans $\mathcal M_3(\C)$.
\item Soient $v_1\in E_1(T)$ et $v_2\in E_2(T)$ non nuls. Trouver $v_3$ tel que $Tv_3=v_2+2v_3$ ; montrer que $(v_1,v_2,v_3)$ est une
base de $\R^3$ et donner une matrice triangulaire semblable à $T$.
\end{enumerate}
\end{exo}

% ---------------------------------------------------------------------
\subsection{Quelle méthode ? Synthèse}\label{ss:synthese-red}

\begin{center}
\renewcommand{\arraystretch}{1.35}
\small
\begin{tabular}{p{0.38\linewidth}p{0.55\linewidth}}
\hline
\textbf{Question} & \textbf{Méthode}\\\hline
$E=F\oplus G$ ? & $F\cap G=\{0\}$ et $\dim F+\dim G=\dim E$ (théorème~\ref{t:adaptee})\\
$F$ stable par $f$ ? & vérifier $f(e_j)\in F$ sur une base de $F$ ; matrice triangulaire par blocs (théorème~\ref{t:blocs})\\
$\lambda$ valeur propre ? & $\rg(A-\lambda I_n)<n$, ou $\det(A-\lambda I_n)=0$ ; $\dim E_\lambda=n-\rg(A-\lambda I_n)$\\
valeurs propres & $\chi_A$ factorisé (méthode~\ref{m:polcar}) ; contrôles par la trace et le déterminant\\
diagonalisable ? & valeurs propres distinctes ; sinon $\dim E_\lambda=m(\lambda)$ pour chaque valeur multiple (méthode~\ref{m:diag})\\
$A^k$, suites couplées & $A=PDP^{-1}$, $A^k=PD^kP^{-1}$ ; coordonnées dans la base propre (méthode~\ref{m:suites})\\
équation matricielle & passer dans la base propre : $N=P^{-1}MP$ ; commutant d'une diagonale (proposition~\ref{p:commute})\\
endomorphisme sur des polynômes, des fonctions & écrire $f(x)=\lambda x$ directement ; chercher des vecteurs propres simples ($X^k$, $e^{\lambda t}$\dots)\\\hline
\end{tabular}
\end{center}

\begin{attention}\textbf{Pièges fréquents.}
\begin{itemize}[leftmargin=1.1em]
\item Un vecteur propre est \emph{non nul} ; $E_\lambda$ contient aussi $0$.
\item Somme directe de trois sous-espaces : les intersections deux à deux ne suffisent pas.
\item Oublier de vérifier que $\chi_A$ est scindé sur $\R$ avant de conclure (rotation).
\item Croire qu'une valeur propre double empêche la diagonalisation : il faut comparer $\dim E_\lambda$ et $m(\lambda)$.
\item Mettre dans $P$ les vecteurs propres en lignes, ou dans un ordre différent de celui de $D$.
\item $A=PDP^{-1}$ et non $P^{-1}DP$, quand les colonnes de $P$ sont des vecteurs propres.
\end{itemize}
\end{attention}

\begin{essentiel}
\begin{itemize}[leftmargin=1.1em,itemsep=2pt]
\item Matrices semblables : même rang, trace, déterminant, polynôme caractéristique.
\item $F$ stable $\iff$ matrice triangulaire par blocs dans une base adaptée ; droite stable $\iff$ vecteur propre.
\item $\chi_A(x)=\det(xI_n-A)$ ; $\Sp(A)$ = racines ; $1\leq\dim E_\lambda\leq m(\lambda)$.
\item Diagonalisable $\iff\sum\dim E_\lambda=n\iff\chi$ scindé et $\dim E_\lambda=m(\lambda)$ ; $n$ valeurs propres distinctes suffisent.
\item $A^k=PD^kP^{-1}$ ; $\tr A=\sum\lambda_i$, $\det A=\prod\lambda_i$ quand $\chi_A$ est scindé.
\end{itemize}
\end{essentiel}

% =====================================================================
%                              CORRIGÉS
% =====================================================================
\clearpage
\section{Corrigés}\label{sec:corriges-red}

\begin{corrige}{red-somme}
\q{1.} Si $v\in F\cap G$, $v=(m,m,m)$ et $3m=0$ : $F\cap G=\{0\}$. De plus $\dim F=2$ (un hyperplan de $\R^3$) et $\dim G=1$ : la
somme est directe et de dimension $3$, donc $\R^3=F\oplus G$. On cherche $v=w+m(1,1,1)$ avec $w\in F$ : la somme des coordonnées
de $w$ est nulle, d'où $x+y+z=3m$. Donc
\[v=\underbrace{\big(x-m,\,y-m,\,z-m\big)}_{\in F}+\underbrace{m(1,1,1)}_{\in G},\qquad m=\frac{x+y+z}3.\]
\q{2.} Les intersections deux à deux sont nulles, mais $(1,0,0)+(0,1,0)-(1,1,0)=0$ est une décomposition non triviale de $0$ :
la somme n'est pas directe (proposition~\ref{p:directe}). D'ailleurs $D_1+D_2+D_3$ est le plan $z=0$, de dimension $2<3$.
\end{corrige}

\begin{corrige}{red-trace}
\q{1.} Par linéarité et $\tr(AB)=\tr(BA)$ : $\tr(AB-BA)=0$, alors que $\tr(I_n)=n\neq0$.
\q{2.} L'endomorphisme $p$ de $\K^n$ associé à $A$ est un projecteur ; dans une base adaptée à $\K^n=\Im p\oplus\Ker p$, sa matrice
est $\mathrm{diag}(I_r,0)$ avec $r=\rg(p)$. Donc $A$ est semblable à cette matrice, et $\tr A=r=\rg A$.
\retenir{la trace ne voit pas l'ordre d'un produit de deux matrices : $\tr(AB-BA)=0$.}
\end{corrige}

\begin{corrige}{red-hyp}
\q{1.} En ajoutant et en retranchant les équations : $x+z=0$ et $y+t=0$. Donc $H_1\cap H_2=\{(x,y,-x,-y)\}$, de base
$\big((1,0,-1,0),(0,1,0,-1)\big)$ : dimension $2=4-2$.
\q{2.} $u\notin H_1$ ($1+1+1+1=4\neq0$) : $\R^4=H_1\oplus\vect(u)$ (proposition~\ref{p:hyperplan}). En revanche $u\in H_2$
($1-1+1-1=0$) : $\vect(u)\subset H_2$, ce n'est pas un supplémentaire.
\end{corrige}

\begin{corrige}{red-stable}
\q{1.} $f(x,0,z)=(3x+z,\,0,\,x+3z)$ a une deuxième coordonnée nulle : $F$ est stable. $f(u)=A(1,1,0)^T=(2,2,0)=2u$ : $G$ est stable
($u$ est un vecteur propre).
\q{2.} $u\notin F$ (sa deuxième coordonnée vaut $1$) : $\R^3=F\oplus G$. Dans la base $(e_1,e_3,u)$ : $f(e_1)=(3,0,1)=3e_1+e_3$,
$f(e_3)=(1,0,3)=e_1+3e_3$, $f(u)=2u$, donc
\[\Mat(f)=\begin{pmatrix}3&1&0\\1&3&0\\0&0&2\end{pmatrix},\qquad\det A=(9-1)\times2=16,\qquad\tr A=3+3+2=8.\]
\retenir{deux sous-espaces stables supplémentaires $\Rightarrow$ matrice diagonale par blocs.}
\end{corrige}

\begin{corrige}{red-elem}
\q{1.} $f$ est linéaire (dérivation, produit par un polynôme fixé) et, si $\deg P\leq n$, $\deg\big((X+1)P'\big)\leq1+(n-1)=n$.
\q{2.} $f\big((X+1)^k\big)=(X+1)\cdot k(X+1)^{k-1}=k(X+1)^k$, et $(X+1)^k\neq0$ : vecteur propre associé à $k$. On obtient $n+1$
valeurs propres distinctes $0,1,\dots,n$ en dimension $n+1$ : ce sont toutes les valeurs propres (corollaire~\ref{c:n-valeurs}),
et chaque sous-espace propre est une droite (théorème~\ref{t:dim-mult}) : $E_k(f)=\vect\big((X+1)^k\big)$. La famille
$\big((X+1)^k\big)_{0\leq k\leq n}$ est une base de vecteurs propres : $f$ est diagonalisable.
\end{corrige}

\begin{corrige}{red-chi}
On développe $\chi_C(x)=\dt{x&0&c\\-1&x&b\\0&-1&x+a}$ selon la première colonne :
\[\chi_C(x)=x\big(x(x+a)+b\big)+\big(0\cdot(x+a)+c\big)=x^3+ax^2+bx+c.\]
Pour les valeurs propres $1,-1,2$, on veut $\chi=(x-1)(x+1)(x-2)=x^3-2x^2-x+2$ : $a=-2$, $b=-1$, $c=2$, soit
$C=\begin{pmatrix}0&0&-2\\1&0&1\\0&1&2\end{pmatrix}$.
\end{corrige}

\begin{corrige}{red-diag1}
\choix{dans $xI_3-A$, chaque ligne a pour somme $x+1$ : on factorise.}
$C_1\leftarrow C_1+C_2+C_3$ puis $L_2\leftarrow L_2-L_1$, $L_3\leftarrow L_3-L_1$ :
\[\chi_A(x)=(x+1)\dt{1&1&1\\1&x-1&1\\1&1&x-1}=(x+1)\dt{1&1&1\\0&x-2&0\\0&0&x-2}=(x+1)(x-2)^2.\]
$A-2I_3$ a trois lignes égales à $(-1,-1,-1)$ : rang $1$, $\dim E_2=2=m(2)$ ; $A$ est diagonalisable. $E_2$ est le plan
$x+y+z=0$, de base $(1,-1,0)$, $(1,0,-1)$ ; $E_{-1}=\vect(1,1,1)$ (somme des lignes de $A$ égale à $-1$). Donc
\[P=\begin{pmatrix}1&1&1\\1&-1&0\\1&0&-1\end{pmatrix},\qquad D=\mathrm{diag}(-1,2,2),\qquad A=PDP^{-1}.\]
Contrôle : $\tr A=3=-1+2+2$, et $A(1,-1,0)^T=(2,-2,0)^T$.
\end{corrige}

\begin{corrige}{red-diag2}
$A_a$ est triangulaire : $\chi_{A_a}=(x-1)^2(x-2)$. La valeur propre $2$ est simple ($\dim E_2=1$). Pour la valeur double $1$ :
$A_a-I_3=\begin{pmatrix}0&a&1\\0&0&1\\0&0&1\end{pmatrix}$. Si $a\neq0$, les deux dernières colonnes sont indépendantes : rang $2$,
$\dim E_1=1<2$, pas diagonalisable. Si $a=0$ : rang $1$, $\dim E_1=2$, diagonalisable. Pour $a=0$ : $E_1=\{z=0\}=\vect(e_1,e_2)$ ;
$(A_0-2I_3)X=0$ donne $x=z$, $y=z$ : $E_2=\vect(1,1,1)$. Ainsi $A_0=PDP^{-1}$ avec
$P=\begin{pmatrix}1&0&1\\0&1&1\\0&0&1\end{pmatrix}$, $D=\mathrm{diag}(1,1,2)$.
\end{corrige}

\begin{corrige}{red-rang1}
\q{1.} Toutes les colonnes de $J_n$ sont égales et non nulles : $\rg J_n=1$. Donc $E_0=\Ker J_n$ est de dimension $n-1$ : c'est
l'hyperplan $x_1+\dots+x_n=0$. Par ailleurs $J_n(1,\dots,1)^T=n(1,\dots,1)^T$ : $n$ est valeur propre. Comme
$\dim E_0+\dim E_n\geq(n-1)+1=n$, il y a égalité, $J_n$ est diagonalisable, $\Sp(J_n)=\{0,n\}$ et $\chi_{J_n}=x^{n-1}(x-n)$.
(Contrôle : $\tr J_n=n=0+\dots+0+n$.)
\q{2.} Si $J_nX=\mu X$, alors $(aI_n+bJ_n)X=(a+b\mu)X$ : les vecteurs propres de $J_n$ sont vecteurs propres de $aI_n+bJ_n$, et
dans la même base, $aI_n+bJ_n$ est diagonale : valeurs propres $a$ (multiplicité $n-1$) et $a+nb$. Donc
$\det(aI_n+bJ_n)=a^{n-1}(a+nb)$, résultat obtenu par opérations élémentaires dans le chapitre \emph{Déterminants}.
\end{corrige}

\begin{corrige}{red-puiss}
$\chi_A=x^2-2x-3=(x-3)(x+1)$, $E_3=\vect(1,1)$, $E_{-1}=\vect(1,-1)$. Avec $P=\begin{pmatrix}1&1\\1&-1\end{pmatrix}$,
$P^{-1}=\frac12\begin{pmatrix}1&1\\1&-1\end{pmatrix}$ :
\[A^k=PD^kP^{-1}=\frac12\begin{pmatrix}3^k+(-1)^k&3^k-(-1)^k\\3^k-(-1)^k&3^k+(-1)^k\end{pmatrix}.\]
Puis $X_k=A^kX_0$ est la première colonne de $A^k$ : $u_k=\frac{3^k+(-1)^k}2$, $v_k=\frac{3^k-(-1)^k}2$.
\end{corrige}

\begin{corrige}{red-racine}
\q{1.} Si $M^2=D$, alors $MD=M\cdot M^2=M^2\cdot M=DM$. Les coefficients de $D$ sont distincts : $M$ est diagonale
(proposition~\ref{p:commute}), $M=\mathrm{diag}(m_1,m_2,m_3)$ avec $m_1^2=1$, $m_2^2=4$, $m_3^2=9$. Les solutions sont les
$8$ matrices $\mathrm{diag}(\pm1,\pm2,\pm3)$.
\q{2.} $A=PDP^{-1}$ ; avec $N=P^{-1}MP$, $M^2=A\iff N^2=D$. L'application $M\mapsto P^{-1}MP$ étant bijective, il y a exactement
$8$ solutions : $M=P\,\mathrm{diag}(\pm1,\pm2,\pm3)\,P^{-1}$.
\end{corrige}

\begin{corrige}{red-nilp}
\q{1.} Voir l'exemple : si $AX=\lambda X$, $X\neq0$, alors $A^pX=\lambda^pX=0$, donc $\lambda=0$. Dans $\C$, $\chi_A$ est scindé et
sa seule racine est $0$ : $\chi_A=x^n$. Par la proposition~\ref{p:trdet}, $\tr A=0$ et $\det A=0$.
\q{2.} Si $A$ est diagonalisable, $A=PDP^{-1}$ avec $D$ diagonale portant les valeurs propres, toutes nulles : $A=0$.
La réciproque est claire. Exemple : $A=\begin{pmatrix}0&1\\0&0\end{pmatrix}$.
\end{corrige}

\begin{corrige}{red-trig}
\q{1.} $\chi_T(x)=x^3-5x^2+8x-4$ (développement, ou contrôle par $\tr T=5$ et $\det T=4$) ; $1$ est racine évidente, et
$\chi_T(x)=(x-1)(x^2-4x+4)=(x-1)(x-2)^2$. Pour la valeur double : $T-2I_3=\begin{pmatrix}3&-17&25\\2&-11&16\\1&-5&7\end{pmatrix}$ est
de rang $2$ (deux premières colonnes non proportionnelles, déterminant nul), donc $\dim E_2=1<2$ : $T$ n'est pas
diagonalisable, ni dans $\R$ ni dans $\C$ (le polynôme et les sous-espaces propres sont les mêmes).
\q{2.} On trouve $E_1=\vect(v_1)$ avec $v_1=(11,7,3)$ et $E_2=\vect(v_2)$ avec $v_2=(3,2,1)$. On cherche $v_3$ tel que
$(T-2I_3)v_3=v_2$ : $v_3=(1,0,0)$ convient ($Te_1=(5,2,1)=2e_1+v_2$). Le déterminant de $(v_1,v_2,v_3)$ vaut
$\dt{11&3&1\\7&2&0\\3&1&0}=7-6=1\neq0$ : c'est une base, et
\[T=Q\begin{pmatrix}1&0&0\\0&2&1\\0&0&2\end{pmatrix}Q^{-1},\qquad Q=\begin{pmatrix}11&3&1\\7&2&0\\3&1&0\end{pmatrix}.\]
\retenir{trigonaliser : vecteurs propres d'abord, puis un vecteur qui complète en donnant une colonne triangulaire.}
\end{corrige}

\end{document}
